% Fig. 1 of the Letter: the two-loop DPS configuration for (a) gluon splittings with effective vertices
% (gg -> HH; gg -> AA with H replaced by A) and (b) massless quark splittings (gg -> ZZ). The red gluon is the
% additional gluon at two loops, drawn in one of its attachments; without it each panel is the one-loop diagram.
% Drawn with TikZ-Feynman; all vertices are placed by hand. Compiled inside dps_diagram.tex by make_figure.py.
\tikzfeynmanset{arrow size=1.2pt}
\colorlet{extra}{red!70!black}                                  % the additional gluon
% labels common to the two panels: daughter momenta, hard subprocesses, panel letter
\newcommand{\commonlabels}[1]{%
  \node at (3.6, 3.12) {\(k\)};
  \node at (3.6,-0.34) {\(\bar r\)};
  \node at (4.66, 1.82) {\(r\)};
  \node at (4.66, 1.00) {\(\bar k\)};
  \node at (5.2, 3.24) {\(\mathcal H_1\)};
  \node at (5.2,-0.44) {\(\mathcal H_2\)};
  \node[extra] at (2.42, 1.4) {\(\ell\)};
  \node[anchor=north west] at (-0.6, 3.75) {#1};}
% the additional gluon, between p_1 - k and p_3 - k
\newcommand{\additionalgluon}{%
  \begin{feynman}
    \vertex (a) at (2.8, 2.1);
    \vertex (b) at (2.8, 0.7);
    \diagram* { (a) -- [gluon, extra] (b) };
  \end{feynman}}
{\small
\begin{tikzpicture}

  % (a) g -> gg with effective vertices: gg -> HH, and gg -> AA with H replaced by A
  \begin{feynman}
    \vertex (g1) at (0.1, 2.8) {\(g(p_1)\)};
    \vertex (g2) at (0.1, 0.0) {\(g(p_2)\)};
    \vertex[dot] (s1) at (2.0, 2.8) {};
    \vertex[dot] (s2) at (2.0, 0.0) {};
    \vertex[crossed dot] (v1) at (5.2, 2.8) {};
    \vertex[crossed dot] (v2) at (5.2, 0.0) {};
    \vertex (x1) at (7.1, 2.8) {\(H(p_3)\)};
    \vertex (x2) at (7.1, 0.0) {\(H(p_4)\)};
    \diagram* {
      (g1) -- [gluon] (s1) -- [gluon] (v1) -- [scalar] (x1),
      (g2) -- [gluon] (s2) -- [gluon] (v2) -- [scalar] (x2),
      (s2) -- [gluon] (v1),                  % p_3 - k: drawn first, it passes under p_1 - k
    };
  \end{feynman}
  \draw[white, line width=6.5pt] (3.46, 1.28) -- (3.74, 1.52);     % gap in p_3 - k at the crossing
  \begin{feynman}
    \diagram* { (s1) -- [gluon] (v2) };      % p_1 - k
  \end{feynman}
  \additionalgluon
  \commonlabels{(a)}

  % (b) g -> q qbar with massless quarks: gg -> ZZ
  \begin{scope}[xshift=8.8cm]
  \begin{feynman}
    \vertex (g1) at (0.1, 2.8) {\(g(p_1)\)};
    \vertex (g2) at (0.1, 0.0) {\(g(p_2)\)};
    \vertex[dot] (s1) at (2.0, 2.8) {};
    \vertex[dot] (s2) at (2.0, 0.0) {};
    \vertex (v1) at (5.2, 2.8);
    \vertex (v2) at (5.2, 0.0);
    \vertex (x1) at (7.1, 2.8) {\(Z(p_3)\)};
    \vertex (x2) at (7.1, 0.0) {\(Z(p_4)\)};
    \diagram* {
      (g1) -- [gluon] (s1) -- [fermion] (v1) -- [boson] (x1),
      (g2) -- [gluon] (s2) -- [fermion] (v2) -- [boson] (x2),
      (v1) -- [plain, with arrow=0.25] (s2), % p_3 - k: drawn first, it passes under p_1 - k
    };
  \end{feynman}
  \draw[white, line width=4pt] (3.45, 1.53) -- (3.75, 1.27);       % gap in p_3 - k at the crossing
  \begin{feynman}
    \diagram* { (v2) -- [plain, with arrow=0.25] (s1) };   % p_1 - k
  \end{feynman}
  \additionalgluon
  \commonlabels{(b)}
  \end{scope}
\end{tikzpicture}}
